\chapter{Checkpoint and Resume}
\label{ch:checkpoint}

\newthought{Long scans get interrupted}---a machine reboots, a job hits its time limit, you press
\code{Ctrl-C}. \Jarvis\ guards against losing that work by checkpointing the sampler's state, so a
run can pick up where it left off.

\section{The checkpoint file}
\label{sec:ckpt-file}

\Jarvis\ keeps a single checkpoint file per sampler, overwritten in place, at:

\begin{jfiletree}
\dirtree{%
.1 \dtfold{<task\_root>/}.
.2 \dtfold{checkpoints/}.
.3 \dtfold{<scan>/}.
.4 \dtfold{<sampler>/}.
.5 \dtfile{state.pkl}\DTcomment{latest checkpoint}.
}
\end{jfiletree}

\noindent Two design choices make this safe:

\begin{description}[style=nextline, leftmargin=1.2em]
  \item[Atomic writes] each save goes to a temporary file and is moved into place with an atomic
        rename, so \file{state.pkl} is never half-written---even if the process dies mid-save.
  \item[Safe barriers] checkpoints are taken on a 30-second heartbeat, but only at points where the
        sampler's state is internally consistent (a ``safe barrier''), never in the middle of an
        update.
\end{description}

\begin{jnote}
With batching (Chapter~\ref{ch:runtime-block}), a safe barrier sits on a \emph{batch boundary}: a
checkpoint is written only once the writer has acknowledged the persisted batch and no batch is still
in flight. That is what lets a resumed run continue without ever duplicating a point.
\end{jnote}

\section{Resuming}
\label{sec:ckpt-resume}

If you re-run a card whose checkpoint exists \emph{without} \cli{-{}-resume}, \Jarvis\ asks what to
do:

\begin{shellbox}
Detected checkpoint file. Re-run from scratch? [y/N] (default: resume in 30s):
\end{shellbox}

\begin{description}[style=nextline, leftmargin=1.2em]
  \item[\code{y} / \code{yes}] start a \emph{fresh} run and discard the existing checkpoint.
  \item[Enter, anything else, or 30\,s timeout] \emph{resume} from the checkpoint.
\end{description}

\noindent To resume non-interactively (in a batch job), pass \cli{-{}-resume}, which skips the
prompt and continues from the latest checkpoint:

\begin{shellbox}
Jarvis bin/my_scan.yaml --resume
\end{shellbox}

\noindent If no checkpoint exists, \Jarvis\ simply starts fresh.

\section{What is and isn't restored}
\label{sec:ckpt-restored}

On resume, \Jarvis\ is deliberate about what it trusts:

\begin{description}[style=nextline, leftmargin=1.2em]
  \item[Configuration comes from the checkpoint] the run continues with the config captured when
        the checkpoint was written, so a resumed run stays consistent with itself.
  \item[The worker factory is rebuilt] the execution machinery is reconstructed fresh from the
        saved blueprint.
  \item[In-flight points are not restored] points that were still running when the run stopped are
        not resurrected; the sampler continues from its last safe barrier.
\end{description}

\section{Inside the checkpoint}
\label{sec:ckpt-payload}

For the curious, the checkpoint is a structured record (not just a pickle blob), with sections
for the run specification, the factory blueprint, the sampler state, and an integrity block:

\begin{description}[style=nextline, leftmargin=1.2em]
  \item[\code{run\_spec}] the normalized config, scan name, task root, workflow modules.
  \item[\code{factory\_blueprint}] worker count, workflow, and module pools to rebuild the factory.
  \item[\code{sampler\_state}] the sampler's signature, random-number state, and (for \textsc{mcmc})
        chains and pending samples.
  \item[\code{integrity}] hashes and signatures used to confirm the checkpoint matches the card.
\end{description}

\begin{jnote}
Because the checkpoint stores the random-number state and the sampler's chains, a resumed
\textsc{mcmc} run continues its chains correctly rather than restarting them---preserving the
statistical validity of the sampling.
\end{jnote}

\section{Summary}
\label{sec:ckpt-summary}

\Jarvis\ checkpoints each sampler to one atomically-written \file{state.pkl} on a 30-second
safe-barrier heartbeat. Re-running prompts you to resume or restart; \cli{-{}-resume} resumes
without asking. Resumed runs use the checkpoint's config, rebuild the factory, and drop in-flight
points.
